\section{Introduction}
\label{sec:introduction}

A fine-tuning corpus is chosen once and used many times. A deployment target is
chosen later, and it is nearly always narrower: a legal question-answering
product does not need the pool's cooking instructions, and the cooking
instructions are not merely useless to it, they compete for capacity and for
optimisation steps. The practical question is which subset of a broad pool to
spend a fixed token budget on when the target distribution is known only
through a handful of examples.

The selection literature answers this question with difficulty. Loss-based
selection keeps the examples the current model predicts worst; gradient-norm
selection keeps the examples that would move the parameters
furthest~\citep{arbogast2021gradnorm}. Both criteria are distribution free, and
that is exactly the problem. An example can be hard because it is informative
about the target, or hard because it is mislabelled, or hard because it belongs
to a region of the pool the target never visits. Under distribution shift the
third case dominates the first, and we measure the consequence in
Section~\ref{sec:analysis}: on our most shifted benchmark, the top decile by
gradient norm contains 61 percent examples from source domains that carry no
target mass at all.

We take the position that selection must be anchored. Given even a small set of
target examples, we can ask a directional question instead of a magnitude
question: does training on this candidate move the parameters in the direction
that a target example would have moved them? That question is answerable with a
sketched gradient inner product, and it separates useful difficulty from
irrelevant difficulty in a way that no magnitude statistic can.

\method{} makes this concrete. It maintains an anchor gradient estimated from a
few hundred target examples, scores every candidate by cosine alignment against
it in a 512-dimensional random projection, and converts those scores into
sampling weights through a temperature that anneals upward during training. The
early curriculum is therefore sharply anchored and the late curriculum
approaches uniform sampling over the pool, which prevents the collapse onto a
narrow slice that a fixed hard threshold produces.

\paragraph{Contributions.}
\begin{itemize}
  \item An anchored selection score computed from sketched per-example
        gradients, with a temperature schedule that turns it into a curriculum
        rather than a static filter (Section~\ref{sec:method}).
  \item An excess-risk bound for the annealed objective that isolates the cost
        of discarded alignment mass from the cost of gradient noise, and a
        corollary giving the anchor size needed for the first term to dominate
        (Section~\ref{sec:method}).
  \item An evaluation on four shifted benchmarks at 1.4B and 7B parameters,
        including the anchor-starvation regime in which \method{} is worse than
        random selection (Sections~\ref{sec:experiments}
        and~\ref{sec:analysis}).
\end{itemize}
